\chapter{Fourier Series}

When Joseph Fourier first made the following claim, the leading
mathematicians of Paris refused to believe it: \emph{any}
periodic signal -- no matter how jagged, how square-cornered, how spiky -- can be built
by adding together pure sine waves whose frequencies are whole-number multiples of one
fundamental frequency. Even a square wave, with its vertical jumps, can be built as a
sum of smooth, rounded sinusoids. A sawtooth ramp is too. The buzz of a violin string,
the vowel in a sung note, the periodic flexing of a tendon as you tap a finger -- all of
them decompose into a stack of sinusoids.

In this chapter the book moves from \emph{describing} signals to
\emph{decomposing} them. In the previous chapters we learned to see a signal as a vector
(Chapter 3) and a sinusoid as a rotating phasor (Chapter 4), and we built an intuition
for what frequency means (Chapter 5). Now we put all three pieces together. The Fourier
series is \emph{the coordinate representation of a periodic signal in a basis made of
sinusoids}. Every idea here is geometry: projection onto basis vectors, orthogonality, and
the Pythagorean theorem -- and each integral is simply a continuous inner product.

\chapterline

\section{The Big Picture: A Periodic Signal Is a Recipe}

Let us fix vocabulary. A signal \(x(t)\) is \emph{periodic} with period \(T\) if
\(x(t+T) = x(t)\) for all \(t\). The smallest such \(T\) is the \emph{fundamental
period}; its reciprocal \(f_0 = 1/T\) is the \emph{fundamental frequency} (in Hz), and
\(\omega_0 = 2\pi/T\) is the fundamental \emph{angular} frequency (in radians per second).
A signal that repeats every \(T\) seconds repeats $f_0$ times per second.

\begin{intuitionbox}
Think of a periodic signal as a single \emph{tile} of length \(T\) stamped over and over
along the time axis. Because the tile repeats forever, all the information lives in one
period. The Fourier series will tell us how to reproduce that one tile by mixing pure
sinusoids -- and once the tile is right, the infinite repetition takes care of itself.
\end{intuitionbox}

The sinusoids we are allowed to use are special: their frequencies must be
\emph{harmonically related}, meaning each is an integer multiple of \(\omega_0\). The
\(k\)-th \emph{harmonic} is \(\cos(k\omega_0 t)\) or \(\sin(k\omega_0 t)\). Why integer
multiples? Because those are exactly the sinusoids that themselves repeat every \(T\)
seconds: if \(\omega = k\omega_0\) then \(\cos(\omega(t+T)) = \cos(k\omega_0 t + 2\pi k)
= \cos(k\omega_0 t)\). A sinusoid at \(2.5\,\omega_0\) would \emph{not} line up with the
tile -- after one period \(T\) it has done two and a half cycles and starts the next tile
out of phase. Only harmonics fit perfectly inside the tile, edge to edge.

\begin{whybox}
\textbf{Harmonics are the building blocks that share the signal's period.} If you add
together any collection of $T$-periodic functions, the sum is still $T$-periodic. So a sum
of harmonics is automatically a candidate for representing a $T$-periodic signal. The central
claim of the Fourier series is the converse: these harmonics form a \emph{complete} family
-- rich enough to build any (reasonable) $T$-periodic signal at all.
\end{whybox}

\subsection{The three forms you will meet}

The Fourier series is written in three equivalent forms. It helps to see all three
at the start so you are never surprised:

\begin{itemize}
\item \textbf{Sine--cosine (trigonometric) form.} A sum
\(\frac{a_0}{2} + \sum_{k\ge1}\bigl(a_k\cos k\omega_0 t + b_k\sin k\omega_0 t\bigr)\).
The most concrete; great for real signals and hand computation.
\item \textbf{Amplitude--phase form.} A sum
\(C_0 + \sum_{k\ge1} C_k\cos(k\omega_0 t - \phi_k)\). Each harmonic is one cosine with its
own amplitude \(C_k\) and phase shift \(\phi_k\). This is how an engineer \emph{hears} it:
``how much of harmonic \(k\), and shifted by how much?''
\item \textbf{Complex-exponential form.} A sum \(\sum_{k=-\infty}^{\infty} c_k\,
e^{\jj k\omega_0 t}\). The cleanest, most symmetric form -- the one that generalizes to the
Fourier transform. The \(c_k\) are complex numbers (rotating phasors). This is the form
we will use most.
\end{itemize}

They carry exactly the same information; we will translate freely between them.

\section{Computing Coefficients Is Projection}

The most important idea in the chapter comes from linear algebra.
In Chapter 3 we learned that the natural inner product for signals on an interval is the
integral of their product. For real $T$-periodic signals we define

\begin{equation}
\inner{x}{y} \;=\; \int_{0}^{T} x(t)\,y(t)\,\dd t,
\label{ch06:inner}
\end{equation}

and over a single period this is a genuine inner product: it is linear, symmetric, and
\(\inner{x}{x} = \int_0^T x(t)^2\,\dd t \ge 0\), zero only for the zero signal. (The
integral can be taken over \emph{any} full period, e.g. \(-T/2\) to \(T/2\); the result is
the same because everything is $T$-periodic.) With an inner product comes geometry:
length, angle, and \emph{orthogonality}.

\begin{connectionbox}
This is the inner product from the linear-algebra companion, made continuous. There,
\(\inner{\vec u}{\vec v} = \sum_i u_i v_i\) sums a product over coordinates. Here we sum
(integrate) a product over time. Time \(t\) plays the role of the index \(i\); the signal
\(x(t)\) is a vector with one component for every instant. Every theorem about
orthonormal bases -- projection, Pythagoras, Parseval -- transfers verbatim.
\end{connectionbox}

\subsection{The harmonics are orthogonal}

Everything works because the harmonic sinusoids are mutually \emph{orthogonal} on one
period: the inner product of two different harmonics is zero. Let us verify the key one.
For integers \(m,n \ge 1\),

\begin{equation}
\int_{0}^{T}\cos(m\omega_0 t)\cos(n\omega_0 t)\,\dd t
= \begin{cases} 0, & m\ne n,\\[2pt] T/2, & m=n. \end{cases}
\label{ch06:ortho-cos}
\end{equation}

\begin{whybox}
The cleanest way to see Eq.~\eqref{ch06:ortho-cos} is geometric, not trigonometric. The
product-to-sum identity gives
\(\cos A\cos B = \tfrac12[\cos(A-B) + \cos(A+B)]\), so the integrand is half a cosine at
frequency \((m-n)\omega_0\) plus half a cosine at \((m+n)\omega_0\). \emph{Integrating a
sinusoid over a whole number of its own periods gives zero} -- you sweep equally through
positive and negative values and they cancel. When \(m\ne n\), both \((m-n)\) and
\((m+n)\) are nonzero integers, so both pieces complete whole cycles over \([0,T]\) and
vanish. When \(m=n\), the first piece is \(\cos(0)=1\), a constant whose integral is
\(T\), giving \(\tfrac12 T\); the second still vanishes. That fact -- ``a sinusoid
integrated over whole cycles averages to zero'' -- is what makes the Fourier
machinery work.
\end{whybox}

By identical arguments,
\(\int_0^T \sin(m\omega_0 t)\sin(n\omega_0 t)\,\dd t = \tfrac{T}{2}\delta_{mn}\) and
\emph{every} sine is orthogonal to \emph{every} cosine:
\(\int_0^T \sin(m\omega_0 t)\cos(n\omega_0 t)\,\dd t = 0\) for all \(m,n\). The sines and
cosines together form an orthogonal basis for $T$-periodic signals -- a coordinate system.

\begin{figure}[H]
\centering
\begin{tikzpicture}
\begin{axis}[
  width=12cm, height=4.2cm,
  xlabel={$t$}, ylabel={},
  domain=0:1, samples=180,
  xmin=0, xmax=1, ymin=-1.25, ymax=1.25,
  axis lines=left, xtick={0,0.5,1}, xticklabels={$0$,$T/2$,$T$},
  ytick={-1,0,1}, enlargelimits=false, legend pos=outer north east,
  legend style={font=\footnotesize}, title style={font=\small}, title={Same harmonic: product stays positive (area $=T/2$)}
]
\addplot[sigblue, thick] {cos(deg(2*pi*x))};
\addlegendentry{$\cos(\omega_0 t)$}
\addplot[specred, thick, dashed] {cos(deg(2*pi*x))};
\addlegendentry{$\cos(\omega_0 t)$}
\path[name path=axis] (axis cs:0,0) -- (axis cs:1,0);
\addplot[black, thick, name path=prod] {cos(deg(2*pi*x))*cos(deg(2*pi*x))};
\addlegendentry{product}
\addplot[fillsig, opacity=0.4] fill between[of=prod and axis];
\end{axis}
\end{tikzpicture}

\begin{tikzpicture}
\begin{axis}[
  width=12cm, height=4.2cm,
  xlabel={$t$}, ylabel={},
  domain=0:1, samples=180,
  xmin=0, xmax=1, ymin=-1.25, ymax=1.25,
  axis lines=left, xtick={0,0.5,1}, xticklabels={$0$,$T/2$,$T$},
  ytick={-1,0,1}, enlargelimits=false, legend pos=outer north east,
  legend style={font=\footnotesize}, title style={font=\small}, title={Different harmonics: product's area cancels to $0$}
]
\addplot[sigblue, thick] {cos(deg(2*pi*x))};
\addlegendentry{$\cos(\omega_0 t)$}
\addplot[specred, thick, dashed] {cos(deg(2*pi*2*x))};
\addlegendentry{$\cos(2\omega_0 t)$}
\path[name path=axis2] (axis cs:0,0) -- (axis cs:1,0);
\addplot[black, thick, name path=prod2] {cos(deg(2*pi*x))*cos(deg(2*pi*2*x))};
\addlegendentry{product}
\addplot[fillsig, opacity=0.4] fill between[of=prod2 and axis2];
\end{axis}
\end{tikzpicture}
\caption{Orthogonality, seen as area. Top: a harmonic times itself is always nonnegative,
so its area over a period is positive ($=T/2$). Bottom: two different harmonics multiply
to a wave whose positive and negative lobes have equal area, so the integral is exactly
zero. This cancellation is why projecting onto one harmonic does not pick up any of the
others.}
\label{ch06:fig-ortho}
\end{figure}

\subsection{Projection gives the coefficients}

Once you have an orthogonal basis, finding coordinates is automatic. If we believe
\(x(t) = \frac{a_0}{2} + \sum_{k\ge1}(a_k\cos k\omega_0 t + b_k\sin k\omega_0 t)\), then
to extract a particular coefficient -- say \(a_n\) -- we take the inner product of both
sides with \(\cos n\omega_0 t\). On the right, orthogonality annihilates every term except
the one we want, which survives with its \(T/2\) self-energy:

\begin{equation}
\inner{x}{\cos n\omega_0 t} = a_n \cdot \frac{T}{2}
\quad\Longrightarrow\quad
a_n = \frac{2}{T}\int_{0}^{T} x(t)\cos(n\omega_0 t)\,\dd t.
\label{ch06:an}
\end{equation}

Identically, \(b_n = \frac{2}{T}\int_0^T x(t)\sin(n\omega_0 t)\,\dd t\), and the constant
term is \(\frac{a_0}{2} = \frac{1}{T}\int_0^T x(t)\,\dd t\), the plain \emph{average} of
the signal (its DC level). Notice we wrote the constant as \(a_0/2\) precisely so that the
formula \eqref{ch06:an} for \(a_n\) works at \(n=0\) too -- a tidy bookkeeping trick.

\begin{happeningbox}
A Fourier coefficient is a \textbf{correlation score}. The integral
\(\int_0^T x(t)\cos(n\omega_0 t)\,\dd t\) asks: ``as I slide along one period, does \(x\)
tend to be large exactly where \(\cos n\omega_0 t\) is large?'' If yes, the product is
mostly positive and the integral is big -- the signal contains a lot of harmonic \(n\).
If \(x\) has no resemblance to that harmonic, the product is positive about as often as
negative and the score is near zero. The Fourier series runs this comparison against
\emph{every} harmonic at once, and the scores are the coordinates. Decomposition is just
``how much does the signal look like each pure tone?''
\end{happeningbox}

\begin{formalbox}
\textbf{Projection formula, abstractly.} For an orthogonal family
\(\{\phi_k\}\) with \(\inner{\phi_k}{\phi_\ell}=\norm{\phi_k}^2\delta_{k\ell}\), the best
representation of \(x\) is \(x = \sum_k \alpha_k \phi_k\) with
\(\alpha_k = \inner{x}{\phi_k}/\norm{\phi_k}^2\). Equation \eqref{ch06:an} is exactly
this with \(\phi_n=\cos n\omega_0 t\) and \(\norm{\phi_n}^2=T/2\). The \(2/T\) is
\(1/\norm{\phi_n}^2\), the basis normalization.
\end{formalbox}

\section{The Complex-Exponential Form}

The sine--cosine form is correct but clumsy: two families of coefficients, special cases at
\(k=0\), and trig identities everywhere. Euler's formula
\(e^{\jj\theta} = \cos\theta + \jj\sin\theta\) lets us fuse the sines and cosines into a
single object. Recall from Chapter 4 that a cosine is the average of two
counter-rotating phasors:

\begin{equation}
\cos(k\omega_0 t) = \tfrac12\bigl(e^{\jj k\omega_0 t} + e^{-\jj k\omega_0 t}\bigr),
\qquad
\sin(k\omega_0 t) = \tfrac1{2\jj}\bigl(e^{\jj k\omega_0 t} - e^{-\jj k\omega_0 t}\bigr).
\label{ch06:euler}
\end{equation}

Substituting these into the trigonometric series and collecting terms by the phasor
\(e^{\jj k\omega_0 t}\) regroups everything into one clean sum that runs over \emph{all}
integers, positive and negative:

\begin{equation}
\boxed{\;x(t) = \sum_{k=-\infty}^{\infty} c_k\, e^{\jj k\omega_0 t},\qquad
c_k = \frac{1}{T}\int_{0}^{T} x(t)\, e^{-\jj k\omega_0 t}\,\dd t.\;}
\label{ch06:complex-fs}
\end{equation}

The synthesis equation (left) \emph{builds} the signal from phasors; the analysis equation
(right) \emph{extracts} each phasor's complex weight. Notice that the analysis integral is
the complex inner product from Chapter~3, $\frac{1}{T}\inner{x}{e^{\jj k\omega_0 t}}$:
conjugating the second slot flips $e^{\jj k\omega_0 t}$ into the counter-rotating
$e^{-\jj k\omega_0 t}$. Multiplying by $e^{-\jj k\omega_0 t}$ freezes the $+k\omega_0$ harmonic into
a stationary arrow whose average survives, while every other harmonic keeps spinning and averages to
zero:
\(\frac1T\int_0^T e^{\jj m\omega_0 t}e^{-\jj n\omega_0 t}\dd t = \delta_{mn}\).

\begin{whybox}
\textbf{Why negative frequencies appear.} A real cosine is built from \emph{two} phasors
spinning in opposite directions; their imaginary parts cancel, leaving a real wave on the
axis. To represent a real signal with single phasors, we therefore need both
\(+k\omega_0\) and \(-k\omega_0\). The negative-frequency term is the
clockwise-spinning partner that keeps the sum real. Concretely, for a real \(x(t)\) the
coefficients come in conjugate pairs: \(c_{-k} = \overline{c_k}\). That single symmetry is
what guarantees the imaginary parts cancel out at every instant (Figure~\ref{ch06:fig-phasor}).
\end{whybox}

\begin{figure}[H]
\centering
\begin{tikzpicture}[scale=1.0]
  \draw[->,gray!60] (-1.6,0) -- (1.6,0) node[right,font=\footnotesize]{Re};
  \draw[->,gray!60] (0,-1.6) -- (0,1.6) node[above,font=\footnotesize]{Im};
  \draw[gray!35] (0,0) circle (1);
  % +k phasor (counterclockwise) at 50 deg
  \draw[->,sigblue,very thick] (0,0) -- (50:1)
        node[midway,above left,font=\footnotesize]{};
  \filldraw[sigblue] (50:1) circle (1.4pt) node[above right,font=\footnotesize]{$c_k e^{\jj k\omega_0 t}$};
  % -k phasor (clockwise) at -50 deg, conjugate
  \draw[->,specred,very thick] (0,0) -- (-50:1);
  \filldraw[specred] (-50:1) circle (1.4pt) node[below right,font=\footnotesize]{$\overline{c_k}\,e^{-\jj k\omega_0 t}$};
  % sum on real axis
  \draw[->,targetgreen,very thick] (0,0) -- ({2*cos(50)/1.55},0);
  \filldraw[targetgreen] ({2*cos(50)/1.55},0) circle (1.4pt);
  \node[targetgreen,below,font=\footnotesize] at ({2*cos(50)/1.55},-0.05) {sum (real)};
  \draw[dashed,gray!55] (50:1) -- ({2*cos(50)/1.55},0);
  \draw[dashed,gray!55] (-50:1) -- ({2*cos(50)/1.55},0);
\end{tikzpicture}
\caption{Why negative frequencies keep a real signal real. The harmonic at $+k\omega_0$
(blue, spinning counterclockwise) and its conjugate partner at $-k\omega_0$ (red, spinning
clockwise) are mirror images across the real axis at every instant. Their imaginary parts
are equal and opposite, so they cancel, and the sum (green) always lands on the real axis.
A real harmonic is this matched pair, never a lone phasor.}
\label{ch06:fig-phasor}
\end{figure}

\subsection{Translating between the forms}

The dictionary is short and worth memorizing. For \(k\ge1\),
\begin{equation}
c_k = \tfrac12(a_k - \jj b_k),\qquad c_{-k} = \tfrac12(a_k + \jj b_k),\qquad c_0 = \tfrac{a_0}{2},
\label{ch06:dict}
\end{equation}
and going the other way, \(a_k = 2\,\mathrm{Re}\,c_k\) and \(b_k = -2\,\mathrm{Im}\,c_k\).
The amplitude--phase form falls right out of the polar form of \(c_k\): writing
\(c_k = \tfrac12 C_k e^{-\jj\phi_k}\), the two conjugate phasors at \(\pm k\omega_0\) recombine into
\(C_k\cos(k\omega_0 t - \phi_k)\). So
\begin{equation}
C_k = 2\abs{c_k} = \sqrt{a_k^2 + b_k^2},\qquad
\phi_k = \arg(a_k + \jj b_k) = \operatorname{atan2}(b_k, a_k).
\label{ch06:ampphase}
\end{equation}
\(C_k\) is how loud harmonic \(k\) is; \(\phi_k\) is where in its cycle it starts. The
\emph{magnitude} spectrum \(\{C_k\}\) is what a spectrum analyzer plots; the \emph{phase}
spectrum \(\{\phi_k\}\) tells you the alignment -- often ignored by ear, but essential for
reconstructing the actual waveform.

\begin{trapbox}
There are several incompatible normalization conventions in the wild. Some books put
\(1/T\) on the synthesis side, some split it as \(1/\sqrt{T}\) on both, some integrate
over \([-\pi,\pi]\). They all describe the same physics but the numbers differ. The
convention here -- \(1/T\) on \emph{analysis} \eqref{ch06:complex-fs}, none on synthesis --
is the most common in engineering and the one that makes \(c_0\) equal the signal average.
When you read a formula elsewhere, \emph{check where the \(1/T\) lives} before trusting a
coefficient value.
\end{trapbox}

\section{Worked Example: The Square Wave}

Let us actually compute something. Take the odd square wave of period \(T\) that equals
\(+1\) on the first half-period and \(-1\) on the second:
\(x(t) = +1\) for \(0 < t < T/2\) and \(x(t) = -1\) for \(T/2 < t < T\). It is the
canonical hard test: perfectly flat plateaus joined by infinitely steep jumps. Can smooth
sinusoids really build those vertical cliffs?

\textbf{Symmetry first -- it does half the work.} The signal is \emph{odd} (\(x(-t) =
-x(t)\) about the midpoint of each plateau). An odd signal has \emph{no} cosine content and
no DC: \(a_k = 0\) for all \(k\), and \(a_0 = 0\). Only sines survive. This is not a
coincidence -- cosines are even, sines are odd, and projecting an odd signal onto an even
basis function gives zero. Symmetry lets you discard half the integrals before lifting a
pen.

\begin{intuitionbox}
Always check symmetry before computing. \textbf{Even} signals (\(x(-t)=x(t)\), mirror-
symmetric) are pure cosine: \(b_k=0\). \textbf{Odd} signals (\(x(-t)=-x(t)\), point-
symmetric) are pure sine: \(a_k=0\). \textbf{Half-wave symmetric} signals (the second half
is the negated first half, as here) have \emph{only odd-numbered harmonics}: \(c_k=0\) for
even \(k\). The square wave is both odd and half-wave symmetric, so only odd sine harmonics
remain. Symmetry is free information.
\end{intuitionbox}

Computing the surviving sine coefficients with \(b_k = \frac2T\int_0^T x(t)\sin(k\omega_0 t)\dd t\),
split into the two half-periods and use \(\omega_0 T = 2\pi\) (so \(k\omega_0 T/2 = k\pi\) and \(\cos(2k\pi)=1\)):
\begin{align}
b_k &= \frac{2}{T}\left[\int_0^{T/2}\!\!\sin(k\omega_0 t)\,\dd t - \int_{T/2}^{T}\!\!\sin(k\omega_0 t)\,\dd t\right]
= \frac{2}{T}\left(\left[-\frac{\cos(k\omega_0 t)}{k\omega_0}\right]_0^{T/2} - \left[-\frac{\cos(k\omega_0 t)}{k\omega_0}\right]_{T/2}^{T}\right) \nonumber\\
&= \frac{2}{k\omega_0 T}\Bigl((1 - \cos k\pi) - (\cos k\pi - 1)\Bigr)
= \frac{2}{k\pi}\bigl(1 - \cos k\pi\bigr)
= \begin{cases} \dfrac{4}{k\pi}, & k \text{ odd},\\[6pt] 0, & k \text{ even}.\end{cases}
\label{ch06:square-bk}
\end{align}
So the square wave's Fourier series is
\begin{equation}
x(t) = \frac{4}{\pi}\left(\sin\omega_0 t + \frac{1}{3}\sin 3\omega_0 t
+ \frac{1}{5}\sin 5\omega_0 t + \frac{1}{7}\sin 7\omega_0 t + \cdots\right)
= \frac{4}{\pi}\sum_{k\,\text{odd}} \frac{\sin k\omega_0 t}{k}.
\label{ch06:square-series}
\end{equation}

The amplitudes fall off like \(1/k\) -- slowly. That slow decay comes from the sharp
edges, and it is why a square wave is hard to reproduce and sounds ``bright''
and buzzy: it is packed with strong high harmonics.

\begin{figure}[H]
\centering
\begin{tikzpicture}
\begin{axis}[
  width=11.5cm, height=5.5cm,
  xlabel={$t$ (one period)}, ylabel={$x(t)$},
  domain=0:1, samples=200,
  xmin=0, xmax=1, ymin=-1.4, ymax=1.4,
  axis lines=left, xtick={0,0.25,0.5,0.75,1},
  xticklabels={$0$,,$T/2$,,$T$}, ytick={-1,0,1},
  enlargelimits=false, legend pos=outer north east, legend style={font=\footnotesize},
  title style={font=\small}, title={Partial sums approaching the square wave}
]
\addplot[gray, very thick, dashed] coordinates {(0,1) (0.5,1)};
\addplot[gray, very thick, dashed] coordinates {(0.5,-1) (1,-1)};
\addlegendentry{ideal square}
\addplot[sourceblue, thick] {(4/pi)*(sin(deg(2*pi*x)))};
\addlegendentry{$N=1$}
\addplot[quotientorange, thick] {(4/pi)*(sin(deg(2*pi*x)) + sin(deg(2*pi*3*x))/3)};
\addlegendentry{$N=3$}
\addplot[targetgreen, thick] {(4/pi)*(sin(deg(2*pi*x)) + sin(deg(2*pi*3*x))/3 + sin(deg(2*pi*5*x))/5 + sin(deg(2*pi*7*x))/7 + sin(deg(2*pi*9*x))/9)};
\addlegendentry{$N=9$}
\end{axis}
\end{tikzpicture}
\caption{Building a square wave from odd harmonics. Even one harmonic captures the gross
shape; adding the third and fifth flattens the plateaus and steepens the edges. As more
harmonics join, the sum hugs the ideal square everywhere except at the jumps -- where an
overshoot remains. That overshoot is the Gibbs phenomenon, next.}
\label{ch06:fig-square}
\end{figure}

\begin{example}[Numbers you can check]
At the midpoint of the first plateau, \(t = T/4\), every odd harmonic
\(\sin(k\omega_0 T/4) = \sin(k\pi/2)\) equals \(+1,-1,+1,\dots\) for \(k=1,3,5,\dots\)
So the series there is \(\frac4\pi(1 - \tfrac13 + \tfrac15 - \tfrac17 + \cdots)\). The
parenthesized sum is the Leibniz series for \(\pi/4\), so the total is
\(\frac4\pi\cdot\frac\pi4 = 1\) -- exactly the height of the plateau. Fourier's formula
also gives us \(\pi/4 = 1 - \tfrac13 + \tfrac15 - \cdots\) as a bonus.
\end{example}

\section{Convergence and the Gibbs Phenomenon}

How can a sum of \emph{continuous} sinusoids equal a \emph{discontinuous} square wave?
The answer is that it does, but not in the naive pointwise sense everywhere. Three
facts settle the matter.

\begin{enumerate}
\item \textbf{At every point where \(x\) is continuous}, the partial sums converge to
\(x(t)\). Smooth stretches are reproduced exactly in the limit.
\item \textbf{At a jump}, the series converges to the \emph{midpoint} of the jump --
the average of the left and right limits. For our square wave at \(t=0\) and \(t=T/2\)
the series gives \(0\), splitting the difference between \(+1\) and \(-1\).
\item \textbf{Near a jump, the partial sums overshoot} by a fixed percentage that never
shrinks. This is the Gibbs phenomenon.
\end{enumerate}

\begin{happeningbox}
The overshoot in Figure \ref{ch06:fig-gibbs} does not go away as you add harmonics -- it
just gets \emph{narrower}. Each finite partial sum overshoots the jump by about
\(8.95\%\) of the jump height (the limiting figure is
\(\frac1\pi\int_0^\pi\frac{\sin u}{u}\dd u - \tfrac12 \approx 0.0895\)), forever. What
\emph{does} shrink is the \emph{width} of the wiggle: it scales like \(1/N\) and squeezes
toward the discontinuity. So the energy of the error goes to zero (the area under the
wiggle vanishes) even though the peak height of the error does not. The series converges in energy,
not in worst-case height. This is exactly why a perfect ``brick-wall'' digital filter rings,
and why an abrupt edit in audio clicks.
\end{happeningbox}

\begin{figure}[H]
\centering
\begin{tikzpicture}
\begin{axis}[
  width=11.5cm, height=5.5cm,
  xlabel={$t$ near the jump at $T/2$}, ylabel={$x(t)$},
  domain=0.35:0.65, samples=200,
  xmin=0.35, xmax=0.65, ymin=-1.45, ymax=1.45,
  axis lines=left, xtick={0.4,0.5,0.6}, xticklabels={,$T/2$,},
  ytick={-1,0,1}, yticklabels={$-1$,$0$,$1$},
  enlargelimits=false, legend pos=outer north east, legend style={font=\footnotesize},
  title style={font=\small}, title={Gibbs overshoot persists at $\approx 9\%$ as $N$ grows}
]
\addplot[gray, very thick, dashed] coordinates {(0.35,1) (0.5,1)};
\addplot[gray, very thick, dashed] coordinates {(0.5,-1) (0.65,-1)};
\addlegendentry{ideal}
\addplot[quotientorange, thick] {(4/pi)*(sin(deg(2*pi*x)) + sin(deg(2*pi*3*x))/3 + sin(deg(2*pi*5*x))/5 + sin(deg(2*pi*7*x))/7 + sin(deg(2*pi*9*x))/9)};
\addlegendentry{$N=9$}
\addplot[sourceblue, thick] {(4/pi)*(sin(deg(2*pi*x)) + sin(deg(2*pi*3*x))/3 + sin(deg(2*pi*5*x))/5 + sin(deg(2*pi*7*x))/7 + sin(deg(2*pi*9*x))/9 + sin(deg(2*pi*11*x))/11 + sin(deg(2*pi*13*x))/13 + sin(deg(2*pi*15*x))/15 + sin(deg(2*pi*17*x))/17 + sin(deg(2*pi*19*x))/19 + sin(deg(2*pi*21*x))/21)};
\addlegendentry{$N=21$}
\addplot[black, dotted, thick] coordinates {(0.35,1.09) (0.5,1.09)};
\node[anchor=south west, font=\tiny, specred] at (axis cs:0.36,1.10) {$1.09$ ($+9\%$)};
\end{axis}
\end{tikzpicture}
\caption{Zoom on the jump. Going from $N=9$ to $N=21$ harmonics squeezes the overshoot
closer to the edge and adds faster ripples, but the peak still reaches roughly $1.09$ --
about $9\%$ past the true value of $1$. More harmonics narrow the error; they do not lower
its peak.}
\label{ch06:fig-gibbs}
\end{figure}

\begin{historybox}
The Fourier series was born in 1807 in one of the most celebrated intellectual controversies in
the history of science. \textbf{Joseph Fourier}, serving as Napoleonic prefect of the Is\`{e}re
department in Grenoble, spent his evenings investigating the diffusion of heat through iron rings
and solid slabs. With only pen, paper, and physical thermometers, Fourier formulated the heat
equation $\partial T/\partial t = \alpha \nabla^2 T$. To solve it for arbitrary initial temperature
profiles $T(x,0) = f(x)$, Fourier proposed that \emph{any} function on an
interval could be expanded as a trigonometric series:
$f(x) = \frac{a_0}{2} + \sum_{n=1}^\infty (a_n\cos nx + b_n\sin nx)$.

When Fourier submitted his 234-page memoir to the Institut de France in Paris on December 21, 1807,
the examining committee -- \textbf{Joseph-Louis Lagrange}, \textbf{Pierre-Simon Laplace},
\textbf{Adrien-Marie Legendre}, and \textbf{Jean-Baptiste Biot} -- reacted with skepticism. Lagrange was especially hostile: decades
earlier he had investigated the vibrating string, concluding that legitimate mathematical functions
could only be represented by power (Taylor) series. Lagrange insisted that an infinite sum of
smooth, infinitely differentiable sinusoidal curves could \emph{never} combine to form a discontinuous
step or an arbitrary curve with sharp corners. The committee rejected Fourier's paper, withholding
the prestigious Grand Prix until 1811, and even then criticized his methods for lack of mathematical
rigor. Fourier was forbidden from publishing his complete work until 1822, when he released his
masterpiece, \emph{Th\'{e}orie analytique de la chaleur}.

Rigor finally arrived in 1829 through \textbf{Peter Gustav Lejeune Dirichlet}. Dirichlet proved the
first exact convergence theorem: if a periodic function is bounded, has a finite number of extrema,
and has a finite number of discontinuities per period (the \emph{Dirichlet conditions}), then the
Fourier series converges to $f(t)$ wherever $f$ is continuous, and converges to the exact arithmetic
midpoint $\frac{1}{2}[f(t^+) + f(t^-)]$ at every jump discontinuity.

One puzzle remained at the jumps. In 1898, American physicist \textbf{Albert A.\ Michelson}
(of the Michelson--Morley experiment) and Samuel Stratton constructed a mechanical
\emph{harmonic analyzer} using 80 calibrated brass springs, gears, and rocking levers to calculate
and draw Fourier sums physically. When they programmed the machine to synthesize a square wave, they
were puzzled: spikes appeared past the top and bottom of every step and did not disappear no
matter how many harmonics were added. Michelson, convinced that the brass springs had an elastic
defect, wrote a perplexed letter to \emph{Nature}. In 1899, Yale mathematical physicist
\textbf{Josiah Willard Gibbs} published the mathematical explanation: the spikes were not a mechanical
error, but an intrinsic mathematical property of truncating the Fourier series of a discontinuous
function. Near any jump of height $H$, the partial sum overshoots by exactly
$\frac{H}{\pi}\int_0^\pi \frac{\sin t}{t}\,\dd t - \frac{H}{2} \approx 0.08949\,H$ -- a persistent
$8.95\%$ overshoot, now known as the \textbf{Gibbs phenomenon} (though first discovered in 1848
by the English mathematician Henry Wilbraham, whose paper had been overlooked).
\end{historybox}

\begin{whybox}
\textbf{When does the series converge nicely, and how fast?} The smoother the signal, the
faster its coefficients die out and the better the convergence. A useful rule of thumb:
if a signal and its first \((p-1)\) derivatives are continuous but the \(p\)-th has jumps,
the coefficients decay like \(1/k^{p+1}\). A square wave (jump in the function itself,
\(p=0\)) decays like \(1/k\). A triangle wave (continuous, but its derivative jumps,
\(p=1\)) decays like \(1/k^2\) -- which is why it has no Gibbs overshoot and converges
visibly faster. A perfectly smooth signal has coefficients that plummet faster than any
power of \(k\). \textbf{Sharp corners cost high harmonics.}
\end{whybox}

\section{Worked Example: The Triangle Wave}

To feel the smoothness-versus-decay rule, compute the triangle wave: the continuous,
even, peak-to-peak \(\pm1\) wave that starts at \(+1\) at \(t=0\), ramps linearly down to \(-1\) at \(t=T/2\), and ramps back up to \(+1\) at \(t=T\). Notice its slope is \(-4/T\) on the first half and \(+4/T\) on the second -- which is \(-4/T\) times our \(\pm 1\) square wave. Integrating term-by-term turns each \(\frac{\sin(k\omega_0 t)}{k}\) in \eqref{ch06:square-series} into \(-\frac{\cos(k\omega_0 t)}{k^2\omega_0}\), dividing every harmonic by a second factor of \(k\) (or you can integrate directly by parts and shift the peak to \(t=0\)). Either way, only odd cosine harmonics survive, falling off as \(1/k^2\):

\begin{equation}
x_{\triangle}(t) = \frac{8}{\pi^2}\left(\cos\omega_0 t - \frac1{9}\cos 3\omega_0 t
+ \frac1{25}\cos 5\omega_0 t - \cdots\right)
= \frac{8}{\pi^2}\sum_{k\,\text{odd}} \frac{(-1)^{(k-1)/2}}{k^2}\cos k\omega_0 t.
\label{ch06:triangle}
\end{equation}

Compare term by term with the square wave \eqref{ch06:square-series}: same odd harmonics,
but \(1/k^2\) instead of \(1/k\). The third harmonic of the square is \(1/3\) of the
fundamental; for the triangle it is only \(1/9\). The high harmonics are much weaker, so just
three terms already trace a near-perfect triangle -- and there is no overshoot, because
the triangle has no jumps for Gibbs to exploit.

\begin{figure}[H]
\centering
\begin{tikzpicture}
\begin{axis}[
  width=11.5cm, height=5.5cm,
  xlabel={$t$ (one period)}, ylabel={$x_\triangle(t)$},
  domain=0:1, samples=200,
  xmin=0, xmax=1, ymin=-1.25, ymax=1.25,
  axis lines=left, xtick={0,0.25,0.5,0.75,1},
  xticklabels={$0$,,$T/2$,,$T$}, ytick={-1,0,1},
  enlargelimits=false, legend pos=outer north east, legend style={font=\footnotesize},
  title style={font=\small}, title={Triangle wave from two harmonics: already nearly exact}
]
\addplot[gray, very thick, dashed] coordinates {(0,1) (0.5,-1) (1,1)};
\addlegendentry{ideal triangle}
\addplot[sourceblue, thick] {(8/pi^2)*(cos(deg(2*pi*x)))};
\addlegendentry{$N=1$}
\addplot[targetgreen, thick] {(8/pi^2)*(cos(deg(2*pi*x)) - cos(deg(2*pi*3*x))/9 + cos(deg(2*pi*5*x))/25)};
\addlegendentry{$N=5$}
\end{axis}
\end{tikzpicture}
\caption{The triangle wave converges far faster than the square. With coefficients dropping
as $1/k^2$, just the fundamental plus two more odd harmonics ($N=5$) hugs the ideal ramp
with no visible overshoot. Compare Figure~\ref{ch06:fig-square}, where the square wave still
shows pronounced ripple at the same harmonic count.}
\label{ch06:fig-triangle}
\end{figure}

\begin{connectionbox}
This is the calculus companion's ``integration is smoothing'' lesson, read in frequency.
The triangle is the running integral of the square; integrating multiplies the \(k\)-th
coefficient by \(1/(\jj k\omega_0)\) -- an extra factor of \(1/k\). Differentiating does
the reverse, multiplying by \(\jj k\omega_0\) and \emph{amplifying} high harmonics. That
is precisely why differentiators are noisy and integrators are calming: the operations
reshape the spectrum by a factor of \(k\).
\end{connectionbox}

\section{The Line Spectrum: What a Periodic Signal Looks Like in Frequency}

Plot the coefficient magnitudes \(\abs{c_k}\) (or \(C_k\)) against frequency \(k\omega_0\)
and you get the \emph{spectrum} of the periodic signal. Because only the harmonic
frequencies \(0, \pm\omega_0, \pm2\omega_0,\dots\) carry any weight, the spectrum is not a
smooth curve -- it is a set of isolated \emph{spikes}, one per harmonic, with nothing in
between. This is a \textbf{line spectrum}, and it is the visual signature of periodicity.

\begin{figure}[H]
\centering
\begin{tikzpicture}
\begin{axis}[
  width=13cm, height=5cm,
  xlabel={harmonic number $k$ (frequency $k\omega_0$)}, ylabel={$\abs{c_k}$},
  xmin=-0.5, xmax=11.5, ymin=0, ymax=0.7,
  axis lines=left, xtick={1,3,5,7,9,11},
  ytick={0,0.2,0.4,0.6}, enlargelimits=false,
  ybar, bar width=4pt, title style={font=\small}, title={Line spectrum of the square wave ($\abs{c_k}=2/(k\pi)$, odd $k$)}
]
\addplot[draw=sigblue, fill=fillsig] coordinates {
  (1,0.6366) (2,0) (3,0.2122) (4,0) (5,0.1273) (6,0)
  (7,0.0909) (8,0) (9,0.0707) (10,0) (11,0.0579)};
\end{axis}
\end{tikzpicture}
\caption{The square wave's one-sided line spectrum. Energy sits \emph{only} at odd
harmonics; even harmonics are exactly absent (half-wave symmetry). The heights fall as
$1/k$ -- the slow decay that makes the wave so rich in highs. A periodic signal lives on a
discrete comb of frequencies, never in between.}
\label{ch06:fig-line}
\end{figure}

\begin{aestheticbox}
The line spectrum is one of the most important pictures in all of signal processing.
\textbf{Periodic in time \(\Leftrightarrow\) discrete (line) in frequency.} A signal that
repeats forever in time uses only a sparse, countable set of frequencies. In the next
chapter we will see the mirror image: a signal that does \emph{not} repeat -- a one-shot
pulse -- spreads its energy across a \emph{continuous} band of frequencies. The Fourier
transform fills in the gaps between the lines as the period \(T\) grows to infinity and
the harmonics \(\omega_0 = 2\pi/T\) crowd together into a continuum. This duality
appears throughout the rest of the book.
\end{aestheticbox}

\begin{applicationbox}
\textbf{Why a violin and a flute playing the same note sound different.} Both produce a
periodic pressure wave at, say, $440$ Hz, so both have a line spectrum on the same comb of
harmonics ($440, 880, 1320, \dots$ Hz). What differs is the \emph{heights} of the lines --
the relative strength of each harmonic. A flute is nearly a pure fundamental with weak
upper harmonics; a violin's bowed string is rich, closer to a sawtooth, with strong highs.
This recipe of harmonic amplitudes is \emph{timbre}. Pitch is which comb; timbre is the
shape of the line spectrum on it. Synthesizers exploit this directly: additive synthesis
literally adds harmonics with chosen $C_k$ to sculpt a tone.
\end{applicationbox}

\section{Energy and Parseval's Theorem}

Because we have a vector resolved along an orthogonal basis, the Pythagorean theorem
applies directly. The average power of the signal over one period equals the sum of the
powers of its individual harmonics. In complex-exponential form this is \emph{Parseval's
theorem for series}:

\begin{equation}
\boxed{\;\frac{1}{T}\int_{0}^{T}\abs{x(t)}^2\,\dd t \;=\; \sum_{k=-\infty}^{\infty}\abs{c_k}^2.\;}
\label{ch06:parseval}
\end{equation}

In trigonometric/amplitude form the same statement reads
\(\frac1T\int_0^T x(t)^2\dd t = \bigl(\tfrac{a_0}{2}\bigr)^2 + \tfrac12\sum_{k\ge1}(a_k^2+b_k^2)
= C_0^2 + \tfrac12\sum_{k\ge1}C_k^2\).

\begin{whybox}
Parseval is just Pythagoras in the function space. The left side is \(\norm{x}^2/T\), the
squared length of the signal vector (its average power). The right side adds up the squared
lengths of its components along each orthogonal basis direction. Because the basis is
orthogonal, there are no cross terms -- the energy of the whole is the sum of the energies
of the parts, exactly as \(\norm{\vec v}^2 = \sum_i v_i^2\) in ordinary geometry. The
harmonics do not ``share'' energy; each owns a private slice, and they add.
\end{whybox}

\begin{example}[Power of the square wave, two ways]
Directly: the square wave is \(\pm1\) everywhere, so \(\frac1T\int_0^T x^2\dd t = 1\). Via
Parseval with \(c_k = 2/(k\pi)\) for odd \(k\) (counting both \(\pm k\)):
\(\sum_k\abs{c_k}^2 = 2\sum_{k\,\text{odd}>0}\bigl(\tfrac{2}{k\pi}\bigr)^2
= \frac{8}{\pi^2}\sum_{k\,\text{odd}}\frac{1}{k^2} = \frac{8}{\pi^2}\cdot\frac{\pi^2}{8} = 1.\)
The two answers agree, and as a side effect we have proved
\(\sum_{k\,\text{odd}}1/k^2 = \pi^2/8\).
\end{example}

\begin{connectionbox}
In Chapter~20 we will compute power spectral density -- where
total power is again distributed across frequency, but now for \emph{random} signals where
the spectrum is continuous and statistical. Parseval is the deterministic ancestor of that
idea: energy conservation across the time/frequency divide. The same accounting governs
both.
\end{connectionbox}

\section{Computing Coefficients in Code}

For signals you only know as samples (the usual case with real sensors), you approximate
the analysis integral \eqref{ch06:complex-fs} by a sum over samples -- which is,
as we will see, exactly the discrete Fourier transform of Chapters 8--9. Here is the
direct, unoptimized version, useful for building intuition before the FFT arrives:

\begin{lstlisting}[language=Python, caption={Fourier coefficients of one period by direct integration.}]
import numpy as np

def fourier_coeffs(x, T, K):
    """Complex FS coefficients c_k for k = -K..K of one period.
    x: samples over [0,T), evenly spaced. Returns dict k -> c_k."""
    N = len(x)
    t = np.arange(N) * (T / N)            # sample times in [0, T)
    w0 = 2 * np.pi / T
    c = {}
    for k in range(-K, K + 1):
        # average of x(t) * exp(-j k w0 t)  ~=  (1/T) * integral
        c[k] = np.mean(x * np.exp(-1j * k * w0 * t))
    return c

# Square wave, one period
N = 1000
x = np.where(np.arange(N) < N // 2, 1.0, -1.0)
c = fourier_coeffs(x, T=1.0, K=7)
print("c_1  =", c[1])         # ~ -0.6366j   -> b_1 = 4/pi
print("|c_3| =", abs(c[3]))   # ~ 0.2122     -> 2/(3 pi)
\end{lstlisting}

\begin{labbox}
Run the snippet, then \emph{resynthesize}: build
\texttt{x\_hat = sum(c[k]*np.exp(1j*k*w0*t) for k in c)} and overlay
\texttt{x\_hat.real} on the original. Watch the Gibbs ears appear at the jumps. Now sweep
\texttt{K} from 1 to 50 and plot the maximum overshoot versus \texttt{K}: it stalls near
$9\%$ and refuses to fall, no matter how high you push \texttt{K}. You have now measured
the Gibbs overshoot yourself.
\end{labbox}

\section{A Note on Real Periodic Signals from Sensors}

So far our square and triangle waves were mathematically perfect. Real periodic signals --
a heartbeat, a footstep cadence, the cyclic motion of a wrist during a repeated gesture --
are only \emph{approximately} periodic, and that gap matters in practice.

\begin{goalbox}
Imagine the wrist-worn controller and a player tapping a steady beat. The accelerometer
sees a roughly periodic bump train. The Fourier series suggests a tempting move: the
fundamental \(\omega_0\) of that bump train \emph{is} the tap rate, and its harmonics
describe the \emph{shape} of each tap (a sharp tap is harmonic-rich; a soft one is nearly
a single sinusoid). So a few low-order coefficients could give us both \emph{tempo} and a
crude \emph{articulation} feature to map to MIDI velocity.

But two limits apply. (1) Real motion is only \emph{quasi}-periodic: the player
drifts in tempo, so a strict series over a fixed \(T\) smears -- we will need the
short-time / windowed view of Chapter 10, not an idealized series. (2) Far more important:
a wrist IMU senses the \emph{gross motion of the whole forearm}, not individual fingers.
No Fourier analysis of an accelerometer can recover \emph{which} finger moved or a
continuous finger-joint angle -- that information is simply not present in the wrist's
acceleration. Detecting individual finger activity needs electromyography (EMG) or force
myography (FMG) over the forearm muscles; continuous finger \emph{pose} needs sensors on
the fingers themselves (e.g. instrumented rings). The Fourier series is a good tool for
analyzing the rhythmic and harmonic \emph{structure} of whatever the sensor actually measures, so
be clear about exactly what that is.
\end{goalbox}

\begin{buildbox}
A first hands-on experiment on real hardware. Sample one axis of an MPU-6050 accelerometer
at, say, $200$ Hz while you tap a finger on a table at a steady rate. Record a few seconds,
then:
\begin{itemize}
\item Estimate the tap period \(T\) (peak-to-peak spacing in the accelerometer trace).
\item Window exactly one period and run the \texttt{fourier\_coeffs} routine above.
\item Inspect \(\abs{c_1}, \abs{c_2}, \abs{c_3}\): a crisp tap shows strong upper
harmonics; a gentle one collapses toward the fundamental.
\end{itemize}
Map \(\omega_0\) to MIDI clock (tempo) and an upper-harmonic ratio like
\(\abs{c_2}/\abs{c_1}\) to note velocity. This is a real, if coarse, gesture-to-MIDI path
that uses only what the IMU genuinely measures. Resist the urge to claim finger identity from it.
\end{buildbox}

\section{Summary}

\begin{recoverybox}
\begin{itemize}
\item A $T$-periodic signal is a sum of harmonically related sinusoids -- frequencies
\(k\omega_0\), \(\omega_0 = 2\pi/T\). These harmonics form an \emph{orthogonal basis}.
\item A Fourier coefficient is a \emph{projection} (inner product) of the signal onto one
harmonic: a correlation score for ``how much of this pure tone is present.''
\item The complex form \(x(t)=\sum_k c_k e^{\jj k\omega_0 t}\),
\(c_k=\frac1T\int_0^T x e^{-\jj k\omega_0 t}\dd t\), is the cleanest; for real signals
\(c_{-k}=\overline{c_k}\).
\item Sharp features (jumps, corners) demand strong high harmonics; smoother signals have
faster-decaying coefficients (\(1/k\) for a square, \(1/k^2\) for a triangle).
\item At a jump the series hits the midpoint and overshoots by \(\approx9\%\) (Gibbs);
the overshoot narrows but never shrinks in height.
\item The spectrum of a periodic signal is a \emph{line spectrum} -- discrete spikes on a
harmonic comb. Periodic in time \(\Leftrightarrow\) discrete in frequency.
\item Parseval: average power equals the sum of squared coefficient magnitudes -- Pythagoras
in function space.
\end{itemize}
\end{recoverybox}

\begin{connectionbox}
\textbf{Where We Are and Where We Go Next}
\begin{itemize}
  \item \textbf{What we have so far}: Any periodic signal decomposes into a countable list of
  harmonic sinusoids at integer multiples of $\omega_0 = 2\pi/T$. Orthogonal projection extracts
  each weight $c_k$, smoothness sets how fast $\abs{c_k}$ falls off ($1/k^{p+1}$), and Parseval's
  theorem equates average time-domain power to the sum of $\abs{c_k}^2$.
  \item \textbf{Where this falls short}:
  \begin{enumerate}
    \item \emph{It requires strict repetition}: The formula assumes $x(t+T) = x(t)$ for all time.
    One-shot events -- a spoken word, a radar pulse, a single wrist flick -- do not repeat every $T$
    seconds.
    \item \emph{Gibbs overshoot at jumps}: Chopping off the series of a waveform with a sharp jump
    leaves an $\approx 9\%$ overshoot spike at every jump, no matter how many harmonics you keep.
  \end{enumerate}
  \item \textbf{Next}: Chapter~7 removes the repetition requirement by letting the period $T \to \infty$.
  As $T$ grows, the harmonic lines pack closer and closer until the discrete sum turns into a
  continuous integral -- the Fourier transform. Chapter~8 then discretizes both time and frequency
  so the whole calculation runs on finite arrays of samples.
\end{itemize}
\end{connectionbox}

\section{Exercises}

\begin{exercise}
For the period-\(T\) signal \(x(t)=A\) for \(0<t<T/4\) and \(0\) elsewhere in
\([0,T)\) (a duty-cycle-\(\tfrac14\) pulse train), compute \(c_0\) and a general
\(c_k\). Show \(c_0\) is the time average \(A/4\), and that \(\abs{c_k}\) follows a
\(\abs{\sinc}\)-like envelope. (This previews the connection between rectangular pulses and
the sinc that dominates the next chapter.)
\end{exercise}

\begin{exercise}
A signal has only the coefficients \(c_0=0\), \(c_{\pm1}=\tfrac12\), \(c_{\pm2}=\tfrac14 \jj\),
all others zero, with \(\omega_0=2\pi\). Write \(x(t)\) as a sum of real sinusoids
(amplitude--phase form). Confirm it is real by checking \(c_{-k}=\overline{c_k}\).
\end{exercise}

\begin{exercise}
Prove the cross-orthogonality \(\int_0^T \sin(m\omega_0 t)\cos(n\omega_0 t)\,\dd t = 0\)
for all integers \(m,n\ge1\). (Hint: use \(\sin A\cos B=\tfrac12[\sin(A+B)+\sin(A-B)]\) and
the whole-cycle averaging argument.)
\end{exercise}

\begin{exercise}
The full-wave rectified cosine \(x(t)=\abs{\cos(\omega_0 t/2)}\) has period \(T=2\pi/\omega_0\)
(since \(\abs{\cos\theta}\) repeats every \(\pi\) radians). Compute its Fourier
coefficients and show that its harmonics decay like \(1/k^2\). Where does its DC
value \(2/\pi\) come from physically? (This is exactly the ripple a rectifier feeds into a
power-supply filter.)
\end{exercise}

\begin{exercise}
Using the square-wave series \eqref{ch06:square-series}, evaluate it at \(t=T/4\) to
re-derive the Leibniz formula \(\pi/4 = 1-\tfrac13+\tfrac15-\cdots\). Then evaluate the
\emph{triangle}-wave series \eqref{ch06:triangle} at its peak and identify the numeric
series you obtain.
\end{exercise}

\begin{exercise}
Sketch (or compute) the line spectrum of \(x(t)=2 + 3\cos(\omega_0 t) - \sin(2\omega_0 t)\).
Give \(c_0, c_{\pm1}, c_{\pm2}\) and the magnitudes \(\abs{c_k}\). Verify Parseval directly
by integrating \(x^2\) over one period and comparing to \(\sum_k\abs{c_k}^2\).
\end{exercise}

\begin{exercise}
A sawtooth wave \(x(t)=t/T\) for \(0\le t<T\) (then repeating) has
\(c_k = \jj/(2\pi k)\) for \(k\ne0\) and \(c_0=\tfrac12\). (a) Verify \(c_0\). (b) Confirm
the \(1/k\) decay and explain, via the smoothness rule, why a sawtooth and a square wave
share the same decay rate. (c) Predict whether the sawtooth shows Gibbs overshoot, and
where.
\end{exercise}

\begin{exercise}
Modify the \texttt{fourier\_coeffs} listing to also \emph{resynthesize} the signal from
\(K\) harmonics and return the maximum reconstruction error. Run it on the square wave for
\(K=5,15,45\). Report the peak overshoot in each case and confirm it approaches roughly
\(9\%\) rather than zero.
\end{exercise}

\begin{exercise}
(Smoothness vs.\ decay.) Construct a continuous, piecewise-linear period-\(T\) signal of
your choice whose first derivative is discontinuous. Predict its coefficient decay rate
from the smoothness rule, compute the first six coefficients numerically, and check the
prediction by fitting \(\log\abs{c_k}\) versus \(\log k\).
\end{exercise}

\begin{exercise}
(Toward the build.) You record a quasi-periodic accelerometer tap train and find the
fundamental is not constant -- the player speeds up over four seconds. Explain in one
paragraph why a single Fourier series over the whole record is the wrong tool, and what you
would do instead (name the chapter-10 technique). Then state plainly what a wrist IMU
\emph{cannot} tell you about individual fingers, and which sensor you would add to get that
information.
\end{exercise}
